\documentclass[12pt]{article}
\usepackage[a4paper,margin=1in]{geometry}
\usepackage{helvet}
\usepackage{courier}
\usepackage{listings}
\usepackage{xcolor}
\usepackage{float}
\usepackage{caption}
\usepackage{subcaption}
\usepackage{graphicx}
\usepackage{hyperref}
\renewcommand{\familydefault}{\sfdefault}
\definecolor{accentblue}{RGB}{0,122,164}
\definecolor{softblue}{RGB}{238,247,253}
\captionsetup{
  labelfont={bf,color=accentblue},
  font=small
}
\newcommand{\evidenceimage}[3]{%
  \begin{figure}[H]
    \centering
    \IfFileExists{#1}{%
      \includegraphics[width=#2\linewidth]{#1}%
    }{%
      \fcolorbox{accentblue!45}{softblue}{%
        \parbox{0.82\linewidth}{\centering\vspace{0.65em}Evidence image pending: \texttt{#1}\vspace{0.65em}}%
      }%
    }%
    \caption{#3}
  \end{figure}
}
\newcommand{\evidencepair}[5]{%
  \begin{figure}[H]
    \centering
    \begin{subfigure}[t]{0.48\linewidth}
      \centering
      \IfFileExists{#1}{\includegraphics[width=\linewidth]{#1}}{%
        \fcolorbox{accentblue!45}{softblue}{\parbox{0.92\linewidth}{\centering\vspace{0.5em}\texttt{#1}\vspace{0.5em}}}%
      }
      \caption{#2}
    \end{subfigure}\hfill
    \begin{subfigure}[t]{0.48\linewidth}
      \centering
      \IfFileExists{#3}{\includegraphics[width=\linewidth]{#3}}{%
        \fcolorbox{accentblue!45}{softblue}{\parbox{0.92\linewidth}{\centering\vspace{0.5em}\texttt{#3}\vspace{0.5em}}}%
      }
      \caption{#4}
    \end{subfigure}
    \caption{#5}
  \end{figure}
}
\lstset{
  basicstyle=\footnotesize\ttfamily,
  breaklines=true,
  frame=single,
  columns=fullflexible,
  showstringspaces=false
}
\title{\textbf{CPT304 Coursework 1: Research-Led Software Enhancement}}
\author{Group XX}
\date{11 May 2026}

\begin{document}
\maketitle

\section{App Selection}
Our group selected BizTrack, an order-management project from the CPT304 project list. It is a small static HTML/CSS/JavaScript app for product, order, and expense tracking, so it was suitable for a research-led enhancement: the business goal is clear, the deployment surface is small, and weaknesses in security, data quality, accessibility, privacy, and testing could be fixed without changing the app's purpose. Our fork is hosted at \url{https://github.com/aaaaxiaoxu/biztrack} and deployed at \url{https://biztrack-6tl.pages.dev/}.

\section{Deficiency 1: DOM XSS Risk}
\subsection{Detection}
Manual source audit found that user-controlled product, order, and expense fields were interpolated into \texttt{innerHTML} in \texttt{products.js}, \texttt{orders.js}, and \texttt{finances.js}. Entering markup such as \texttt{<img onerror=alert(1)>} could be parsed as HTML instead of displayed as data.
\subsection{Literature}
OWASP explains that output encoding prevents XSS by ensuring untrusted input is treated as data rather than executable browser code \cite{owasp-xss}. The DOM-based guidance also recommends using safe DOM sinks such as text nodes where possible.
\subsection{Implementation}
The fix follows OWASP's safe-sink principle: rows are now built with \texttt{createElement()} and values are assigned using \texttt{textContent}. Action icons became real buttons with event listeners, avoiding inline string handlers that mixed data and code.
\begin{lstlisting}[caption={Before and after: safe table output}]
// Before
prodRow.innerHTML = `<td>${product.prodName}</td>`;

// After
const cell = document.createElement("td");
cell.textContent = product.prodName;
prodRow.appendChild(cell);
\end{lstlisting}

\section{Deficiency 2: CSV Formula Injection}
\subsection{Detection}
The export functions joined raw values with commas. If a product name began with \texttt{=}, \texttt{+}, \texttt{-}, or \texttt{@}, Excel or LibreOffice could interpret it as a formula when the CSV was opened.
\subsection{Literature}
OWASP defines CSV Injection as embedding untrusted input into exported spreadsheet cells; cells beginning with formula metacharacters may execute spreadsheet functions \cite{owasp-csv}. Recommended mitigation includes quoting fields and prefixing dangerous values.
\subsection{Implementation}
A shared \texttt{BizTrackCore.generateCSV()} now escapes quotes, wraps comma/newline fields, and prefixes formula-like values with an apostrophe. All three exports use this single tested helper.
\begin{lstlisting}[caption={Before and after: CSV export}]
// Before
Object.values(row).join(",");

// After
function escapeCsvValue(value) {
  const text = String(value ?? "");
  const safe = /^[=+\-@\t\r]/.test(text) ? `'${text}` : text;
  return /[",\n]/.test(safe) ? `"${safe.replace(/"/g, '""')}"` : safe;
}
\end{lstlisting}

\section{Deficiency 3: Keyboard and Focus Accessibility}
\subsection{Detection}
Manual keyboard testing found invisible focus indicators because the global CSS removed outlines. Sortable table headers and icon-only controls relied on pointer clicks and lacked clear keyboard/reader semantics.
\subsection{Literature}
WCAG 2.2 requires keyboard operability and visible focus indicators for interactive content \cite{wcag22}. W3C's Focus Visible note explains that removing visible focus harms users who navigate without a mouse \cite{wai-focus}.
\subsection{Implementation}
The global outline removal was replaced with \texttt{:focus-visible}. Dynamic edit/delete controls are now \texttt{button} elements with \texttt{aria-label}; sortable headers receive keyboard activation. This turns previously mouse-only actions into accessible controls.
\begin{lstlisting}[caption={Before and after: actionable controls}]
// Before
<i onclick="deleteProduct('PD001')" class="delete-icon"></i>

// After
deleteButton.type = "button";
deleteButton.setAttribute("aria-label", `Delete product ${prodID}`);
deleteButton.addEventListener("click", () => deleteProduct(prodID));
\end{lstlisting}

\section{Deficiency 4: Fragile Local Data and Numeric Validation}
\subsection{Detection}
The app parsed \texttt{localStorage} directly with \texttt{JSON.parse()} and accepted parsed numbers without checking for \texttt{NaN}. A corrupted storage entry could stop page rendering; invalid amounts could poison totals.
\subsection{Literature}
MDN describes Web Storage values as strings that must be parsed by application code \cite{mdn-storage}. web.dev also warns that local storage is synchronous and should be handled carefully with clear fallbacks \cite{web-storage}.
\subsection{Implementation}
The new core layer wraps storage parsing in \texttt{try/catch}, falls back to default data, validates non-negative numeric fields, and generates transaction IDs from the maximum existing ID rather than array length.
\begin{lstlisting}[caption={Before and after: resilient storage}]
// Before
orders = JSON.parse(localStorage.getItem("bizTrackOrders"));

// After
orders = core.parseStoredArray(localStorage, "bizTrackOrders", defaultOrders);
orderTotal = core.calculateOrderTotal(itemPrice, qtyBought, shipping, taxes);
\end{lstlisting}

\section{Baseline Standards Evidence}
This section records the five mandatory baseline standards required by the brief. The evidence combines deployment logs, automated test output, Lighthouse audit results, bilingual interface screenshots, and privacy-compliance screenshots. Each item is shown with a numbered figure and a short caption so that the marker can verify the standard directly.

\subsection{Live Uptime}
BizTrack is deployed through Cloudflare Pages from the \texttt{main} branch. The deployment evidence shows a successful production build and publish process, including repository cloning, application build, asset upload, and deployment to Cloudflare's global network. This confirms that the submitted URL is a live production deployment rather than a local-only prototype.

\evidenceimage{Figures/部署log.png}{0.92}{Production deployment evidence for BizTrack on Cloudflare Pages.}

\subsection{Test Coverage}
The project uses Vitest with V8/Istanbul coverage. The latest local run completed successfully with 4 test files and 15 tests passing. The reported coverage is 95.86\% statements, 87.17\% branches, 98.38\% functions, and 98.49\% lines, which exceeds the 80\% minimum baseline.

\evidenceimage{Figures/coverage.png}{0.92}{Automated test coverage evidence from Vitest/Istanbul.}

\subsection{Accessibility}
The Lighthouse audit reports an Accessibility score of 100. This result reflects the refactoring work on keyboard-focus visibility, labelled controls, semantic buttons, table interaction, and navigation structure. Manual checks are still useful, but the automated score demonstrates that the implementation meets the coursework threshold of 90+.

\evidenceimage{Figures/英文lighthouse.png}{0.82}{Lighthouse accessibility evidence showing a score of 100.}

\subsection{Bilingual Function (i18n)}
BizTrack includes a visible language selector in the header. Switching the selector updates the interface from English to Chinese, including navigation labels, page titles, dashboard headings, form text, and translated UI content. The language preference is persisted in browser storage, so the selected locale remains active across page reloads. Figures below show the same application area in English and Chinese, proving that the bilingual toggle is functional rather than a static screenshot.

\evidencepair{Figures/英文.png}{English interface with the language selector set to EN.}{Figures/中文.png}{Chinese interface after switching the language selector.}{Internationalization evidence showing BizTrack in English and Chinese.}

\subsection{Legal Compliance}
The application provides a dedicated Privacy Policy page explaining that business records, language preference, and cookie-consent preference are stored locally in the browser. It also clarifies that the project does not use accounts, analytics trackers, advertising cookies, or server-side personal-data collection. Together with the cookie-consent banner, this satisfies the legal-compliance baseline for a responsible client-side application.

\evidenceimage{Figures/隐私政策.png}{0.92}{Privacy Policy page evidence for legal-compliance baseline.}

\section{Individual Contribution Forms}
Attach the completed peer-assessment spreadsheet for all four members. Add direct pull-request links for each member here before submission.

\section{References}
\begin{thebibliography}{9}
\bibitem{owasp-xss} OWASP Foundation, ``Cross Site Scripting Prevention Cheat Sheet,'' OWASP Cheat Sheet Series. [Online]. Available: \url{https://cheatsheetseries.owasp.org/cheatsheets/Cross_Site_Scripting_Prevention_Cheat_Sheet.html}
\bibitem{owasp-csv} OWASP Foundation, ``CSV Injection,'' OWASP. [Online]. Available: \url{https://owasp.org/www-community/attacks/CSV_Injection}
\bibitem{wcag22} W3C, ``Web Content Accessibility Guidelines (WCAG) 2.2,'' W3C Recommendation. [Online]. Available: \url{https://www.w3.org/TR/WCAG22/}
\bibitem{wai-focus} W3C WAI, ``Understanding Success Criterion 2.4.7: Focus Visible.'' [Online]. Available: \url{https://www.w3.org/WAI/WCAG21/Understanding/focus-visible}
\bibitem{mdn-storage} MDN Web Docs, ``Using the Web Storage API.'' [Online]. Available: \url{https://developer.mozilla.org/en-US/docs/Web/API/Web_Storage_API/Using_the_Web_Storage_API}
\bibitem{web-storage} web.dev, ``Storage for the web.'' [Online]. Available: \url{https://web.dev/storage-for-the-web/}
\end{thebibliography}
\end{document}
